\bookchapter{Capítulo 41: Panorama das ferramentas de programação com IA em 2026}

\chapterintro{Aviso antecipado: este capítulo é uma “fotografia” do ecossistema de agosto de 2026. Em seis meses estará desatualizado — e essa é a primeira lição. O objetivo não é dizer “use esta ferramenta”, e sim construir um critério de escolha. Quando você entende a filosofia de cada ferramenta, decide bem em qualquer época.}

\booksection{Três escolas e um quadro de decisão}

As ferramentas de 2026 se dividem em três grandes escolas, cada uma com uma filosofia própria. Veja o panorama:

\begin{center}
\begin{tabularx}{\linewidth}{|X|X|X|X|}
\hline
\textbf{Escola} & \textbf{Ferramentas} & \textbf{Filosofia} & \textbf{Cenário} \tabularnewline
\hline
IA integrada à IDE & Cursor 3, Trae, Windsurf, CodeBuddy & Integrar a IA profundamente ao editor & Desenvolvimento diário, code review, melhoria incremental \tabularnewline
Agente de terminal & Claude Code, Codex CLI, Aider & Interface no terminal; a IA opera arquivos e comandos diretamente & Operações em escopo de projeto, mudanças em lote, CI/CD \tabularnewline
Builder na nuvem & v0, Bolt, Replit Agent & Criar projetos completos do zero no navegador & Prototipagem rápida, início sem configuração, pré-visualização imediata \tabularnewline
\hline
\end{tabularx}
\end{center}

\begin{center}\small\color{gray}As três escolas das ferramentas de IA\end{center}

A questão não é “qual é a melhor”, e sim “qual combina com o seu cenário”. Três critérios úteis:

\begin{itemize}
  \item Qual é a sua base? Iniciante Builder na nuvem; base sólida IA na IDE; avançado agente de terminal.
  \item Em que fase está o projeto? Do zero Builder; iteração IA na IDE; manutenção/refatoração agente de terminal.
  \item Qual é a necessidade de deploy? Validação rápida Builder; projeto sério IA na IDE; automação agente de terminal.
\end{itemize}

\booksection{Escola da IA integrada à IDE}

É a escola mais comum: embute a IA no editor, preservando a experiência tradicional e acrescentando reforços. Em 2026, os destaques são Cursor 3, Trae, Windsurf e CodeBuddy.

\booksubsection{Cursor 3}

Cursor é a IDE de IA mais usada no mundo. O Cursor 3, de abril de 2026, marcou uma virada: o modo Agent ultrapassou o autocomplete por Tab — de “a IA completa o código” para “a IA conclui a tarefa”. Recursos centrais: edição em múltiplos arquivos, percepção de contexto (estrutura e dependências do projeto) e integração profunda com o terminal (a IA roda comandos e lê a saída dentro do editor).

\booksubsection{Trae}

Trae é a IDE de IA da ByteDance, com grande base de usuários na China. O modo SOLO (ou Builder) é o clássico “criar um projeto com uma frase”: você descreve a necessidade e a IA faz o caminho todo, da inicialização ao código. Seus pontos fortes são o suporte em chinês e a adaptação à rede local — ideal para começar rápido.

\booksubsection{Outras ferramentas}

O Windsurf (antigo Codeium) é famoso pela experiência fluida de autocomplete. O CodeBuddy da Tencent Cloud, como extensão do VS Code, integra-se bem ao ecossistema de nuvem e a projetos de WeChat Mini Programs. Alibaba Tongyi Lingma e Baidu Comate também oferecem integração semelhante, cada um com otimizações próprias para o cenário chinês.

\booksection{Escola do agente de terminal}

É a escola que mais cresceu em 2026. Ideia central: a IA não deve ser só um plugin do editor; deve operar o ambiente de desenvolvimento como um dev — ler e escrever arquivos, rodar comandos, ver saídas, instalar dependências. O representante é o Claude Code.

\booksubsection{Claude Code}

Claude Code é a ferramenta de terminal da Anthropic e, em meados de 2026, alcançou US\$ 2,5 bilhões de receita anualizada. Funciona de outro jeito: você descreve a tarefa em linguagem natural no terminal e ele lê os arquivos, altera o código, roda testes, analisa resultados e corrige em ciclo — sem abrir o editor.

A vantagem é a “percepção do projeto inteiro”: sem a limitação da janela do editor, ele navega pelo diretório e entende as dependências. Também suporta o CLAUDE.md, para que a IA já saiba o contexto no início da sessão. Em refatoração, mudanças em lote e escrita de testes em projetos grandes, o agente de terminal costuma ser mais eficiente que plugins de IDE.

\booksubsection{Codex CLI e Aider}

O Codex CLI da OpenAI oferece experiência semelhante no terminal e suporta /import, trazendo bases de código existentes para a IA entender. O Aider é open source, aceita vários modelos (GPT, Claude, DeepSeek etc.) e é bom para quem gosta de trocar de modelo. A filosofia é a mesma: dar à IA controle pleno do ambiente de desenvolvimento no terminal.

\booksection{Escola do Builder na nuvem}

É a mais indicada para iniciantes. Sem instalar nada: abra o navegador, descreva a necessidade e a IA cria o projeto completo na nuvem, com pré-visualização imediata.

\begin{itemize}
  \item v0 (Vercel): ótimo para componentes e páginas de UI, protótipos web rápidos.
  \item Bolt: cria aplicações full stack a partir de uma frase, com front, back e banco.
  \item Replit Agent: trabalha na IDE na nuvem do Replit, cria, roda e publica o projeto.
\end{itemize}

Vantagem comum: “começar sem configuração”. Limite: com o projeto crescendo, o ambiente na nuvem fica mais rígido que a IDE local. São ideais para a fase “de 0 a 1”: validar a ideia rápido e, depois, migrar para a IDE local e desenvolver a fundo.

\booksection{Ferramentas mudam; metodologia, não}

Se amanhã todas as ferramentas de programação com IA desaparecessem, três hábitos deveriam permanecer: pensar antes de agir, dar um pequeno passo por vez e verificar ao terminar. Ferramentas são amplificadores; a metodologia é a fundação. Amplificadores podem ser trocados, a fundação não pode balançar.

Olhando 2024–2026: em 2024 o GitHub Copilot era o rei; em 2025, o Cursor o superou; em 2026, o Claude Code despontou. A cada ferramenta nova, alguém diz “dessa vez é diferente”, mas a conclusão é a mesma: quem se dá bem já tem metodologia sólida. Essas pessoas se adaptam em poucos dias; quem só decorou cliques recomeça do zero a cada troca.

Por isso, a trilha sugerida é: começar com uma ferramenta (qualquer uma) e praticar as cinco regras; com a metodologia firme, experimentar outras escolas e sentir as diferenças; por fim, montar a sua combinação — por exemplo, Cursor no dia a dia, Claude Code em refatoração e testes, v0 em protótipos. A combinação é pessoal; a metodologia é comum.

\begin{tcbpractice}{Mão na massa: monte o seu "arsenal"}
Posicione-se pelos três critérios (base, fase do projeto, deploy): iniciante experimente Bolt ou v0 com uma página; com base Cursor ou Trae em uma tarefa de iteração; para se desafiar rode o Claude Code no terminal. Depois responda: qual escola combina com a sua fase atual?

\end{tcbpractice}

\begin{tcbsummary}{Resumo do capítulo}
\begin{itemize}
  \item Três escolas: IA na IDE (dia a dia), agente de terminal (projeto inteiro), Builder na nuvem (protótipo para iniciantes)
  \item Três perguntas: base de programação? fase do projeto? necessidade de deploy?
  \item No Cursor 3, o modo Agent superou o Tab; o Claude Code chegou a US\$ 2,5 bi anuais
  \item O Builder na nuvem serve para "de 0 a 1"; depois, migre para a IDE local
  \item Ferramentas mudam a cada três meses: quem tem método se adapta em poucos dias
\end{itemize}

\end{tcbsummary}
